\section{Same-image Learned-representation Audit}\label{app:learned-audit}

\subsection{Scope, representations and provenance}
This retrospective extension asks whether prototype similarity, name recognition
and centered painter agreement describe the same response in a common learned
representation. These comparisons connect the prompted-artist evaluation of
\citet{su2025prompted} with CLIP \citep{radford2021} and the CSD style descriptors
of \citet{somepalli2024style}. The protocol was fixed after the 31-feature analysis
and before inspecting learned-representation outcomes. No new images, painters,
prompts or reference works enter this extension.

The input census contains all 1,878 originals: 1,008 generated images, 649 primary
references and 221 development works. An additional 131 views apply the previously
recorded painting-region crops to 90 references and 41 development images, giving
2,009 extraction rows per representation. The region sensitivity substitutes these
views for their originals; unchanged works retain their original view. It does not
double-count cropped works or alter generated pixels.

Both encoders return 768-dimensional unit vectors. We use their Euclidean/cosine
geometry without coordinate IQR scaling, whitening or selected dimensions.
Table~\ref{tab:learned-provenance} identifies the fixed public checkpoints and
local extraction implementations. CSD uses its released style projection of the
CLIP visual backbone, followed by unit normalization; the content head is unused.
Its release documentation reports a discrepancy between released weights and
published benchmark numbers. Our CSD measurements therefore characterize the
specified public checkpoint, without claiming to reproduce those benchmark numbers.

\begin{table}[htbp]
\caption{Learned-extraction provenance. Revision and SHA-256 prefixes are shortened
to 12 characters here; retained records contain complete identifiers.}
\label{tab:learned-provenance}
\centering\small\setlength{\tabcolsep}{6pt}
\begin{tabularx}{\linewidth}{@{}lXXX@{}}
\toprule
Encoder & Checkpoint revision & Weight SHA-256 & Implementation SHA-256\\
\midrule
CLIP ViT-L/14 & \texttt{32bd64288804} & \texttt{a2bf730a0c7d} & \texttt{156750d6d709}\\
CSD ViT-L/14 & \texttt{5bc26a6fb048} & \texttt{40e92fad63a3} & \texttt{787b95697b6c}\\
\bottomrule
\end{tabularx}
\end{table}

The checkpoint repositories are \texttt{openai/clip-vit-large-patch14} and
\texttt{tomg-group-umd/CSD-ViT-L}. The retained OpenAI vision implementation is
from commit \texttt{a1d071733d71}; CSD architecture and preprocessing follow
commit \texttt{3a9df32605b8}. Strict tensor loading checks the selected model
state and style-head dimensions. The CSD adapter has its own versioned input
record; completed CLIP inputs, code and receipts are preserved.

\subsection{Pixels, native crops and numerical checks}
Each source is hash-verified. EXIF orientation and valid ICC-to-sRGB conversion
precede the optional audited region. The resulting image undergoes bicubic resize
to a 224-pixel short side, a $224\times224$ center crop, and CLIP channel
normalization, directly from original pixels. The primary learned view is thus a
native center crop, unlike the full-frame hand-feature measurement.

The Hugging Face CLIP processor uses floor center offsets. Native CSD follows
torchvision's rounded offsets; these can differ by one pixel when the resized
long side is odd. A separately bound CSD adapter implements that convention,
fixed before CSD outcomes. The difference is retained rather than silently
changing completed CLIP extraction. Processed float32 tensors receive individual
hashes. Models run serially in evaluation mode, float32, with MPS batches of four.

The supplementary checks use four deterministically selected cases: a
generated image, an original reference, an odd-aspect reference, and an audited
region. CPU/MPS and batch-four/single comparisons have maximum absolute
unit-vector-coordinate difference $8.11\times10^{-7}$ for CLIP and
$1.26\times10^{-6}$ for CSD, below the declared
$5\times10^{-5}$ tolerance; their tensor hashes match extraction. This checks
backend and batching consistency on those cases, not all possible images or
artistic validity. Supplemental provenance binds configurations, implementation
files and the dependency lock. The CLIP environment records Python 3.13.11,
PyTorch 2.13.0, Transformers 4.57.6, NumPy 2.5.2 and Pillow 11.3.0. Numerical
replay uses retained embeddings and hash-bound row identities; exact inference
can vary across environments.

\subsection{Prototype gain and labeled agreement}
For $A=4$ painters, let $g_a$ be the mean of named unit embeddings across scenes
and repeats, $g_0$ the corresponding generic-painting mean, and $\mu_a$ the mean
reference embedding. These prototypes are \emph{unnormalized} means, so
$g_a^\top\mu_b$ equals average image-to-reference cosine similarity. Define
$\bar g=A^{-1}\sum_a g_a$, $\bar\mu=A^{-1}\sum_a\mu_a$,
$d_a=g_a-\bar g$, $r_a=\mu_a-\bar\mu$ and $H=\sum_a\|r_a\|^2$.
The equally painter-weighted diagonal gain has the exact identity
\begin{equation}
\underbrace{\frac1A\sum_a g_a^\top\mu_a-g_0^\top\bar\mu}_{\text{prototype gain}}
=\underbrace{(\bar g-g_0)^\top\bar\mu}_{\text{common contribution}}
+\underbrace{\frac1A\sum_a d_a^\top r_a}_{\text{labeled contribution}}
=(\bar g-g_0)^\top\bar\mu+\frac{H}{A}\widehat\beta.
\label{eq:learned-prototype-gain}
\end{equation}
The same decomposition holds for the artist-free baseline. A positive prototype
gain can therefore include movement common to every painter name. Recognition
is a separate calculation: each generated image is assigned to the largest dot
product with a \emph{unit-normalized} reference prototype. We retain confusion
matrices and macro-average development-work recognition over painters.

All 24 reference-label permutations are descriptive controls. If $P_\pi$ is the
permuted diagonal prototype score, then
\begin{equation}
P_\pi=\bar g^\top\bar\mu+\frac{H}{A}\widehat\beta_\pi,
\qquad
\widehat D_\pi=1+\widehat Q-2\widehat\beta_\pi.
\label{eq:learned-permutation}
\end{equation}
Here $H$ and $\widehat Q$ are unchanged by reference-label permutation. Thus
ranking permutations by higher $P_\pi$, higher $\widehat\beta_\pi$ or lower
$\widehat D_\pi$ gives the same ordering by construction. These ranks are
neither independent corroboration nor confirmatory permutation tests.

\subsection{Sensitivity scope and artifact access}
The audit retains all configurations and painter pairs, repeats the original
contrast/error/calibration calculations in each representation, and reports
common/specific changes from both control arms. Reference contrast energy is
specific to the representation and view; normalized errors across them are not
a shared perceptual scale. Leave-one-scene-out ranges describe influence, not
confidence intervals. The 221-work development target is an alternative finite
panel already used in this project, not a fresh collection or service replication.

CSD derives from CLIP, and both are applied to the same images. Their agreement
is therefore evidence about sensitivity to related representations. All four
painters occur in CSD's published style-tag list, and training-image overlap is
unknown for both encoders. Neither representation supplies perceptual ground truth,
an unseen-artist test or validation of the AI-coded source regions.

A deterministic local inventory verifies the hashes, byte counts and detected
media types of exactly 1,878 originals; a companion attribution inventory binds
870 historical sources to retained license, credit and source-page metadata.
Recorded reproduction-file licenses comprise 763 public-domain entries, 23 CC0
entries and 84 CC BY/CC BY-SA entries. These records distinguish the complete
cohort from the four displayed public-domain examples. The inventories make local
artifact contents reviewable without copying pixels. They do not constitute a
public image release or establish complete live recovery; the public exact-pixel
access limitation remains.
Richer retained metadata matches each acquisition by hash and supplies license
URLs and credit fields for all 84 attribution-bearing files. Two lack a separate
artist-text field; release preparation requires contextual review of these credits.
